\chapter{ทบทวนความรู้พื้นฐานเรื่องการแปรผัน}
\section{การแปรผัน}
การแปรผัน แบ่งออกเป็น 3 ชนิด ดังนี้
\begin{enumerate}
\item การแปรผันตรง
\item การแปรผันแบบผกผัน
\item การแปรผันเกี่ยวเนื่อง
\end{enumerate} 

\subsection{การแปรผันตรง}
\textbf{การแปรผันตรง} คือความสัมพันธ์ระหว่างปริมาณสองปริมาณ โดยที่อัตราส่วนของปริมาณหนึ่งต่ออีกปริมาณหนึ่งเป็นค่าคงตัว กล่าวคือ เมื่อปริมาณหนึ่งเปลี่ยนแปลง อีกปริมาณหนึ่งจะเปลี่ยนแปลงไปในทิศทางเดียวกันและเป็นสัดส่วนกัน  

\begin{definitionbox}\label{definition1.1}
กำหนดให้ $x$ และ $y$ แทนปริมาณใด ๆ  
จะกล่าวว่า $y$ \textbf{เป็นสัดส่วนโดยตรง} กับ $x^n$ เมื่อ $n$ เป็นจำนวนจริง  
เขียนแทนด้วยสัญลักษณ์ $y \propto x^n$  
ซึ่งหมายความว่า
\[
y = \alpha x^n
\]
โดยที่ $\alpha$ เป็นค่าคงตัว และ $\alpha \neq 0$ เรียก $\alpha$ ว่า \textbf{ค่าคงตัวของการแปรผัน}\\
\medskip
\textbf{กรณีพิเศษ} เมื่อ $n=1$ จะเรียกว่า $y$ \textbf{แปรผันตรง} กับ $x$
\end{definitionbox}

\newpage
\subsection{การแปรผกผัน}
\textbf{การแปรผกผัน} คือความสัมพันธ์ระหว่างปริมาณสองปริมาณ โดยที่ผลคูณของปริมาณทั้งสองเป็นค่าคงตัว กล่าวคือ เมื่อปริมาณหนึ่งเพิ่มขึ้น อีกปริมาณหนึ่งจะลดลงในอัตราส่วนที่สอดคล้องกัน  

\begin{definitionbox}\label{definition1.2}
กำหนดให้ $x$ และ $y$ แทนปริมาณใด ๆ ที่ไม่เป็นศูนย์  
จะกล่าวว่า $y$ \textbf{แปรผกผัน} กับ $x$ เขียนแทนด้วยสัญลักษณ์ $y \propto \dfrac{1}{x}$  
เมื่อสามารถเขียนความสัมพันธ์ในรูป
\[
y = \frac{\alpha}{x}
\]
โดยที่ $\alpha$ เป็นค่าคงตัวและ $\alpha \neq 0$ เรียก $\alpha$ ว่า \textbf{ค่าคงตัวของการแปรผัน}
\end{definitionbox}

\subsection{การแปรผันเกี่ยวเนื่อง}
\textbf{การแปรผันเกี่ยวเนื่อง} คือความสัมพันธ์ที่ปริมาณหนึ่งขึ้นอยู่กับหลายปริมาณพร้อมกัน ซึ่งอาจเป็นการแปรผันตรง การแปรผกผัน หรือการผสมกันของทั้งสองลักษณะ 

\begin{definitionbox}\label{definition1.3}
กำหนดให้ $x_1, x_2, x_3, \ldots, x_n$ และ $y$ แทนปริมาณใด ๆ  
จะกล่าวว่า $y$ \textbf{แปรผันเกี่ยวเนื่อง} กับ $x_1, x_2, \ldots, x_n$  
เมื่อ $y$ แปรผันตรงกับผลคูณของปริมาณเหล่านั้น เขียนได้เป็น
\[
y = \alpha x_1 x_2 \cdots x_n
\]
โดยที่ $\alpha$ เป็นค่าคงตัวและ $\alpha \neq 0$ เรียก $\alpha$ ว่า \textbf{ค่าคงตัวของการแปรผัน}
\end{definitionbox}

\newpage
\begin{examplebox}[NETSAT คณิตศาสตร์ กุมภาพันธ์ 2565]  
กำหนดให้ $y$ เป็นสัดส่วนโดยตรงกับ $x^2$ และเมื่อ $x=-2$ จะได้ $y=-5$  
ข้อใดต่อไปนี้ถูกต้อง
\begin{multicols}{2} \begin{enumerate}
\item เมื่อ $x=2$ จะได้ $y=-5$
\item เมื่อ $x=-\dfrac{1}{2}$ จะได้ $y=-\dfrac{5}{6}$
\item เมื่อ $x=\dfrac{1}{2}$ จะได้ $y=-\dfrac{5}{6}$
\item ถูกต้องทั้งข้อ 1, 2 และ 3
\end{enumerate} 
\end{multicols}
\end{examplebox}
\newpage
\begin{examplebox}[NETSAT คณิตศาสตร์ กุมภาพันธ์ 2567]  
กำหนดให้ $y$ เป็นสัดส่วนโดยตรงกับ $x^2$ และเมื่อ $x=-6$ จะได้ $y=9$  
ถ้า $x=8$ แล้ว $y$ มีค่าเท่าใด
\begin{multicols}{2} \begin{enumerate}
\item $16$
\item $32$
\item $64$
\item $256$
\end{enumerate} \end{multicols}
\end{examplebox}
\newpage
\begin{examplebox}[NETSAT คณิตศาสตร์ สิงหาคม 2565]  
กำหนดให้ $y$ เป็นสัดส่วนโดยตรงกับ $x^2$ และเมื่อ $x=-1$ จะได้ $y=3$  
ข้อใดเป็นสมการที่แสดงความสัมพันธ์ระหว่าง $y$ และ $x$
\begin{multicols}{2} \begin{enumerate}
\item $y=x^2+2$
\item $y=-x^2+4$
\item $y=2x^2+1$
\item $y=3x^2$
\end{enumerate} \end{multicols}
\end{examplebox}
\newpage
\begin{examplebox}[NETSAT คณิตศาสตร์ 2566]
กำหนดให้ $y$ เป็นสัดส่วนโดยตรงกับ $x^2$ และเมื่อ $x=4$ จะได้ $y=-4$  
ข้อใดต่อไปนี้เป็นสมการที่แสดงความสัมพันธ์ระหว่าง $y$ และ $x$
\begin{multicols}{2} \begin{enumerate}
\item $y=\dfrac{1}{2} x^2-12$
\item $y=-\dfrac{1}{2} x^2+4$
\item $y=\dfrac{1}{2} x^2-8$
\item $y=-\dfrac{1}{4} x^2$
\end{enumerate} \end{multicols}
\end{examplebox}
